\documentclass[runningheads]{llncs}
\usepackage[T1]{fontenc}
\usepackage{graphicx}
\usepackage{booktabs}
\usepackage[utf8]{inputenc}
\usepackage{amsmath}
\usepackage{listings}
\usepackage{xcolor}
\usepackage{amssymb}
\usepackage{cite}
\usepackage{hyperref}
\usepackage{subcaption}
\usepackage{xspace}

\definecolor{codegray}{rgb}{0.95,0.95,0.95}
\definecolor{codegreen}{rgb}{0,0.6,0}
\definecolor{codepurple}{rgb}{0.58,0,0.82}
\definecolor{backcolour}{rgb}{0.95,0.95,0.92}

\begin{document}

\title{Review of Restaking Protocol Security:\\
Threat Models, Attack Surfaces, and Mitigation Strategies\\
in EigenLayer and Emerging Shared-Security Protocols}

\author{Author information scrubbed for double-blind reviewing}

\maketitle

\begin{abstract}
Restaking protocols have emerged as a transformative innovation in blockchain security, enabling staked assets to secure multiple networks and services simultaneously. EigenLayer, the pioneering restaking protocol, has grown to over \$15B in total value locked (TVL) by 2024, with emerging competitors including Symbiotic, Karak, and Babylon. This review analyzes security threats, attack surfaces, and mitigation strategies across restaking protocols through comprehensive synthesis of available academic literature, security audits, and protocol documentation. We identify five primary threat categories: (1) dual-layer slashing risks from Ethereum and AVS penalties, (2) smart contract vulnerabilities in restaking and AVS contracts, (3) centralization risks from validator concentration, (4) liquidity and unstaking challenges with LST/LRT tokens, and (5) cryptoeconomic attacks including correlated failures and AVS collusion. Our analysis reveals that while restaking enhances capital efficiency and shared security, it introduces novel systemic risks including correlated slashing across multiple AVSs, validator centralization pressures, and complex interdependencies between LST issuers, operators, and AVSs. We synthesize mitigation strategies including slashing caps, risk diversification mechanisms, formal verification approaches, and governance frameworks. Critical research gaps include formal game-theoretic models of restaking security, empirical analysis of slashing incidents, and composability risks with DeFi protocols. This review provides actionable security guidelines for protocol designers, validators, and delegators operating in the restaking ecosystem.
\end{abstract}

\keywords{Restaking protocols, EigenLayer, shared security, slashing risks, cryptoeconomic security, security review}

\vspace{-0.3cm}
\section{Introduction}

Restaking protocols represent a paradigm shift in blockchain security economics~\cite{eigenlayer2024,eigenlayerwhitepaper2023}, enabling staked assets to secure multiple networks and services beyond their native blockchain. This innovation addresses a fundamental limitation in proof-of-stake (PoS) systems~\cite{buterin2014ethereum,wood2014ethereum,buterin2017casper}: the underutilization of staked capital that remains locked while providing security to only a single network. By allowing validators to restake their staked ETH or liquid staking tokens (LSTs)~\cite{lido2024,rocketpool2024,coinbase2024cbeth} to secure additional services called Actively Validated Services (AVSs), restaking protocols create a shared security model~\cite{wood2021polkadot} that enhances capital efficiency while potentially strengthening the security of emerging protocols.

EigenLayer, launched in mid-2023, pioneered the restaking concept on Ethereum and has achieved remarkable adoption, with total value locked (TVL) growing from \$1.1B in January 2024 to over \$20B by June 2024, representing a 1,335\% increase~\cite{defillama2024}. By late 2024, EigenLayer had become one of the largest DeFi protocols by TVL, demonstrating the significant demand for restaking services. The protocol enables Ethereum validators and LST holders to restake their assets to secure AVSs, which can include data availability layers, oracle networks, bridges, and other middleware services. This success has spawned a competitive ecosystem of restaking protocols, including Symbiotic~\cite{symbiotic2024,symbioticwhitepaper2024} (supporting diverse ERC-20 tokens), Karak~\cite{karak2024} (multi-asset and cross-chain restaking), and Babylon~\cite{babylon2024} (Bitcoin-native restaking), each with distinct architectural choices and security models.

However, the rapid growth and complexity of restaking protocols introduce significant security challenges that have not been systematically analyzed. Restaking creates novel attack surfaces including: (1) \textbf{Dual-Layer Slashing:} Validators face slashing penalties from both Ethereum consensus and individual AVSs, creating compounded risk exposure, (2) \textbf{Correlated Failures:} A single validator misbehavior can trigger slashing across multiple AVSs simultaneously, (3) \textbf{Centralization Pressures:} Economic incentives may concentrate validator power, creating systemic risks, (4) \textbf{Smart Contract Complexity:} The integration of restaking contracts with LST protocols and AVSs introduces new vulnerability vectors, and (5) \textbf{Liquidity Risks:} Unstaking delays and LST/LRT (Liquid Restaking Token) depegging can create cascading failures.

The security implications of restaking are profound. Unlike traditional staking where risks are isolated to a single network, restaking creates interconnected risk networks where failures can propagate across multiple services. A validator slashing event could simultaneously impact Ethereum consensus, multiple AVSs, and associated LST/LRT protocols, potentially creating systemic failures affecting billions of dollars in staked capital.

Despite the critical importance of restaking security, no comprehensive review exists that synthesizes threats, attack surfaces, and mitigation strategies across the restaking ecosystem. Existing analyses are fragmented across security audit reports~\cite{certora2025,code4rena2023,trailofbits2024,openzeppelin2024}, technical blog posts~\cite{blocksec2024,sigmaprime2024,cobo2024,okx2024}, and protocol documentation~\cite{eigenlayerwhitepaper2023,symbioticwhitepaper2024}, with limited academic treatment of restaking security as a distinct research domain~\cite{durvasula2024robust,zhang2024proof,li2024mitigating}.

\subsection{Research Questions}

This review addresses the following research questions:

\begin{enumerate}
    \item \textbf{RQ1 (Architecture):} How do current restaking protocols architect their shared-security model, including AVS design, operator sets, slashing conditions, and LST/LRT support mechanisms?
    
    \item \textbf{RQ2 (Threats):} What are the main security threats and failure modes introduced by restaking protocols, including smart contract vulnerabilities, correlated slashing risks, economic attacks, AVS collusion, and governance risks?
    
    \item \textbf{RQ3 (Mitigations):} Which technical and cryptoeconomic mitigation mechanisms are proposed or deployed in restaking protocols, including slashing rules, risk caps, diversification strategies, insurance mechanisms, and formal verification approaches?
    
    \item \textbf{RQ4 (Risk Distribution):} How are risks allocated among validators, delegators, AVSs, LST issuers, and end users in different restaking protocol designs, and what are the implications for security and incentive alignment?
    
    \item \textbf{RQ5 (Research Gaps):} Which aspects of restaking security are underexplored in current literature, including formal game-theoretic models, empirical incident analysis, composability risks with DeFi protocols, and long-term sustainability?
\end{enumerate}

\subsection{Contributions}

The main contributions of this review are:

\begin{itemize}
    \item \textbf{Comprehensive Threat Taxonomy:} A systematic classification of security threats across five primary categories (dual-layer slashing, smart contract vulnerabilities, centralization risks, liquidity challenges, cryptoeconomic attacks) based on analysis of protocol documentation, security audits, and academic literature.
    
    \item \textbf{Protocol Comparison Framework:} A comparative analysis of four major restaking protocols (EigenLayer, Symbiotic, Karak, Babylon) examining architectural differences, security models, and risk profiles.
    
    \item \textbf{Mitigation Strategy Synthesis:} A comprehensive synthesis of proposed and deployed mitigation mechanisms, including technical solutions (slashing caps, formal verification) and cryptoeconomic mechanisms (risk diversification, insurance, governance).
    
    \item \textbf{Risk Distribution Analysis:} A systematic analysis of how risks are allocated across different stakeholders (validators, delegators, AVSs, LST issuers) and the security implications of these allocations.
    
    \item \textbf{Research Gap Identification:} Identification of critical underexplored areas including formal security models, empirical analysis, and composability risks.
    
    \item \textbf{Actionable Security Guidelines:} Evidence-based recommendations for protocol designers, validators, delegators, and AVS operators to enhance restaking security.
\end{itemize}

\subsection{Paper Organization}

The remainder of this paper is organized as follows: Section~\ref{sec:related} reviews related work on staking security, shared security models, and cryptoeconomic mechanisms. Section~\ref{sec:background} provides background on restaking protocols and their security models. Section~\ref{sec:methodology} describes our review methodology. Section~\ref{sec:architecture} analyzes restaking protocol architectures. Section~\ref{sec:threats} presents our threat taxonomy. Section~\ref{sec:mitigations} reviews mitigation strategies. Section~\ref{sec:risk_distribution} analyzes risk distribution across stakeholders. Section~\ref{sec:comparison} provides comparative analysis of major protocols. Section~\ref{sec:gaps} identifies research gaps. Section~\ref{sec:discussion} discusses findings and implications. Section~\ref{sec:conclusion} concludes.

\vspace{-0.3cm}
\section{Related Work}
\label{sec:related}

This section positions our review within existing literature on staking security, shared security models, and cryptoeconomic mechanisms.

\subsection{Proof-of-Stake Security}

Extensive research has examined security properties of proof-of-stake consensus mechanisms~\cite{buterin2014ethereum,wood2014ethereum,buterin2017casper,kiayias2017ouroboros}. Key security properties include liveness (chain progress), safety (no double-spending), and economic security (cost of attacks)~\cite{li2022survey,zhou2022state,liu2021survey}. Slashing mechanisms penalize validators for misbehavior, including double-signing, downtime, and protocol violations~\cite{atzei2017survey}. Research has analyzed slashing in various PoS systems including Ouroboros~\cite{kiayias2017ouroboros}, Tendermint~\cite{buchman2016tendermint,kwon2014tendermint}, and Ethereum 2.0~\cite{buterin2017casper}. However, traditional PoS security analysis assumes isolated networks without cross-protocol dependencies introduced by restaking. Recent work by Durvasula and Roughgarden~\cite{durvasula2024robust} provides initial game-theoretic analysis of restaking security, characterizing robust security as a function of attack costs and profits.

\subsection{Liquid Staking Protocols}

Liquid staking protocols enable stakers to receive liquid tokens (LSTs) representing their staked assets, enabling participation in DeFi while maintaining staking rewards~\cite{zhang2022cross,li2021sovereign}. Major LST protocols include Lido (stETH)~\cite{lido2024}, Rocket Pool (rETH)~\cite{rocketpool2024}, and Coinbase (cbETH)~\cite{coinbase2024cbeth}, collectively representing over \$50B in staked ETH~\cite{ethereum2024staking,defillama2024}. Research has analyzed LST security including slashing risks, validator centralization concerns~\cite{chen2022rollup}, and depegging vulnerabilities~\cite{li2021sovereign,trailofbits2023staking}. However, restaking introduces additional complexity by enabling LSTs to secure multiple services simultaneously, creating novel interdependencies. Liquid Restaking Tokens (LRTs) such as Renzo~\cite{renzo2024}, EtherFi~\cite{etherfi2024}, Puffer~\cite{puffer2024}, and Kelp DAO~\cite{kelp2024} further complicate the security landscape by adding another layer of tokenization and protocol dependencies.

\subsection{Shared Security Models}

Shared security models enable multiple protocols to leverage a common validator set, enhancing security for smaller networks~\cite{wang2023survey,belchior2021survey}. Research has examined shared security in contexts including Polkadot parachains~\cite{wood2021polkadot}, Cosmos zones~\cite{kwon2019cosmos}, and Ethereum Layer 2 solutions~\cite{kalodner2018arbitrum,buterin2022danksharding}. Cross-chain interoperability research~\cite{herlihy2018atomic,zhang2021cross} provides additional context for shared security models. However, restaking represents a distinct model where security is shared across diverse services (AVSs) rather than blockchain networks, introducing different threat vectors. Li~\cite{li2024mitigating} examines how EigenLayer and Lido address Ethereum PoS challenges, while Zhang and Wang~\cite{zhang2024proof} propose verification mechanisms applicable to AVSs. Das et al.~\cite{das2025reliablocks} developed Reliablocks during the EigenLayer hackathon, contributing to AVS reliability measurement.

\subsection{Cryptoeconomic Security}

Cryptoeconomic security relies on economic incentives and penalties to ensure protocol security~\cite{zhou2020high,zhou2021high,chen2021sok,qin2021atomic}. Research has analyzed slashing mechanisms, validator economics, and attack costs in PoS systems~\cite{kiayias2017ouroboros,buterin2017casper}. Flash loan attacks and MEV extraction have been studied in DeFi contexts~\cite{gao2021flashloan,zhou2020high,zhou2021high,chen2021sok}, but restaking introduces multi-layer slashing where validators face penalties from both base layer and AVSs, creating complex risk interactions. Durvasula and Roughgarden~\cite{durvasula2024robust} provide initial game-theoretic analysis, but comprehensive models are still lacking. Smart contract security research~\cite{huang2021survey,chen2020defining,ahmed2021survey,luu2016making,rodler2018sereum} provides foundations, but restaking-specific vulnerabilities require additional analysis.

\subsection{Research Gap}

While individual aspects of restaking security have been discussed in audit reports~\cite{certora2025,code4rena2023,trailofbits2024,openzeppelin2024,consensys2024restaking} and technical analyses~\cite{blocksec2024,sigmaprime2024,cobo2024,okx2024}, and recent academic work has begun exploring restaking~\cite{durvasula2024robust,zhang2024proof,li2024mitigating,das2025reliablocks}, no systematic review exists that synthesizes threats, attack surfaces, and mitigation strategies across the restaking ecosystem. Related work on blockchain security~\cite{li2022survey,zhou2022state}, smart contract security~\cite{huang2021survey,chen2020defining,atzei2017survey}, cross-chain interoperability~\cite{zhang2022cross,li2021sovereign,wang2023survey,belchior2021survey,li2023survey}, and DeFi security~\cite{chen2021sok,gao2021flashloan} provides context, but restaking-specific analysis is limited. Our work addresses this gap by providing comprehensive analysis through structured synthesis of available literature, synthesizing findings from 57 sources including academic papers, security audits, and protocol documentation.

\vspace{-0.3cm}
\section{Background on Restaking Protocols}
\label{sec:background}

This section provides fundamental concepts necessary to understand restaking protocol security.

\subsection{Restaking Fundamentals}

Restaking enables staked assets to secure multiple services beyond their native blockchain. The core innovation is allowing validators and LST holders to commit their staked capital to secure AVSs, which can include:

\begin{itemize}
    \item \textbf{Data Availability Layers:} Services providing data availability for Layer 2 solutions
    \item \textbf{Oracle Networks:} Price feeds and off-chain data providers
    \item \textbf{Bridge Protocols:} Cross-chain asset transfer services
    \item \textbf{Middleware Services:} Various infrastructure services requiring validation
\end{itemize}

Restaking creates a shared security model where AVSs leverage Ethereum's validator set without requiring separate token issuance or validator recruitment.

\subsection{Major Restaking Protocols}

\subsubsection{EigenLayer}

EigenLayer, launched in 2023, enables native ETH stakers and LST holders to restake their assets. Key features include:
\begin{itemize}
    \item Support for native ETH restaking and LST restaking (stETH, rETH, cbETH)
    \item Operator-based architecture where operators run AVS software
    \item Delegation model allowing restakers to delegate to operators
    \item Slashing mechanisms for AVS misbehavior
    \item TVL exceeding \$15B by late 2024
\end{itemize}

\subsubsection{Symbiotic}

Symbiotic, launched in 2024, differentiates by supporting diverse ERC-20 tokens beyond ETH/LSTs:
\begin{itemize}
    \item Universal staking framework supporting multiple token types
    \item Flexible AVS integration model
    \item \$29M Series A funding in 2025
\end{itemize}

\subsubsection{Karak}

Karak offers multi-asset and cross-chain restaking:
\begin{itemize}
    \item Support for ETH, LSTs, LRTs, WBTC, stablecoins, Pendle PTs
    \item Cross-chain operation across Ethereum, Arbitrum, BSC, Blast, Mantle, K2
    \item Enhanced liquidity through diverse asset support
\end{itemize}

\subsubsection{Babylon}

Babylon enables Bitcoin-native restaking:
\begin{itemize}
    \item Direct BTC staking through self-custodial vaults
    \item No wrapping or bridging required
    \item Extends Bitcoin utility to secure other blockchains
\end{itemize}

\subsection{Security Model Components}

Restaking security depends on several components:

\begin{itemize}
    \item \textbf{Slashing Mechanisms:} Penalties for validator/operator misbehavior
    \item \textbf{Operator Sets:} Validators running AVS software
    \item \textbf{Delegation:} Mechanism for restakers to delegate to operators
    \item \textbf{Risk Caps:} Limits on restaking exposure per validator
    \item \textbf{AVS Governance:} Decision-making for slashing conditions and upgrades
\end{itemize}

\vspace{-0.3cm}
\section{Review Methodology}
\label{sec:methodology}

This systematic review follows PRISMA (Preferred Reporting Items for Systematic Reviews and Meta-Analyses) guidelines to ensure comprehensive and reproducible results.

\subsection{Review Protocol}

We developed a review protocol defining: (1) research questions, (2) search strategy, (3) inclusion/exclusion criteria, (4) quality assessment, and (5) data extraction procedures.

\subsection{Search Strategy}

We conducted searches across multiple sources to identify relevant literature:

\textbf{Academic Databases:}
\begin{itemize}
    \item IEEE Xplore
    \item ACM Digital Library
    \item SpringerLink
    \item ScienceDirect
    \item arXiv (cs.CR, cs.CY, q-fin.CP)
    \item Google Scholar
\end{itemize}

\textbf{Gray Literature:}
\begin{itemize}
    \item Security audit reports (Certora, Trail of Bits, Code4rena, OpenZeppelin)
    \item Security research blogs (BlockSec, Sigma Prime, Cobo, OKX Research)
    \item Protocol documentation and whitepapers
    \item GitHub repositories
\end{itemize}

\textbf{Search Terms:}
\begin{itemize}
    \item "restaking" OR "re-staking"
    \item "EigenLayer" OR "shared security"
    \item "liquid restaking" OR "LRT"
    \item "actively validated service" OR "AVS"
    \item "slashing risk" AND ("restaking" OR "EigenLayer")
    \item "restaking security" OR "restaking vulnerability"
\end{itemize}

\textbf{Date Range:} 2023-2025 (restaking is a very recent innovation)

\textbf{Search Execution:} Conducted in January 2025

\subsection{Inclusion and Exclusion Criteria}

\textbf{Inclusion Criteria:}
\begin{itemize}
    \item Papers/reports discussing security, threats, or risks of restaking protocols
    \item Security audits of restaking protocols or AVSs
    \item Technical analyses of restaking mechanisms and vulnerabilities
    \item Protocol documentation describing security models
    \item Published between 2023-2025
    \item Written in English
\end{itemize}

\textbf{Exclusion Criteria:}
\begin{itemize}
    \item Papers not focusing on security aspects
    \item Duplicate publications
    \item Insufficient technical detail
    \item Papers focusing solely on performance without security implications
    \item Marketing materials without technical content
\end{itemize}

\subsection{Study Selection Process}

Given the novelty of restaking protocols (launched 2023-2024), the literature base is still emerging. Our review synthesizes available academic work, security audits, and protocol documentation through a structured approach:

\begin{enumerate}
    \item \textbf{Academic Literature Search:} We searched academic databases including IEEE Xplore, ACM Digital Library, SpringerLink, ScienceDirect, and arXiv using keywords: "restaking", "EigenLayer", "shared security", "liquid restaking", "AVS", and related terms. We identified peer-reviewed papers, preprints, and technical reports addressing restaking security.
    
    \item \textbf{Gray Literature Collection:} We collected security audit reports from major auditing firms (Certora, Code4rena, Trail of Bits, OpenZeppelin), technical analyses from security research teams (BlockSec, Sigma Prime, Cobo, OKX Research), and official protocol documentation from EigenLayer, Symbiotic, Karak, and Babylon.
    
    \item \textbf{Inclusion Criteria:} Sources were included if they: (1) addressed security, risks, or vulnerabilities in restaking protocols, (2) provided technical analysis or empirical data, (3) were published or updated between 2023-2025, and (4) were accessible in English.
    
    \item \textbf{Quality Assessment:} We prioritized peer-reviewed academic papers, published security audits from reputable firms, and official protocol documentation. Technical blog posts and analyses were included when they provided unique insights not available in academic literature.
\end{enumerate}

Our final included sources comprise academic papers (including foundational work on PoS security and recent restaking-specific research), security audit reports, technical analyses, and protocol documentation. The field's rapid evolution means new sources continue to emerge, and this review captures the state of knowledge as of early 2025.

Figure~\ref{fig:prisma} presents a conceptual flow diagram illustrating our study selection process. Given the field's novelty and limited peer-reviewed literature, we include both academic sources and high-quality gray literature (security audits, technical analyses, protocol documentation) with appropriate quality assessment.

\begin{figure}[!ht]
    \centering
    \includegraphics[width=0.9\textwidth]{figures/prisma.pdf}
    \caption{Conceptual flow diagram showing the study selection process for this review. The diagram illustrates the multi-stage approach from initial search across academic databases and gray literature to final inclusion based on quality and relevance criteria. Given the field's novelty, we synthesize both peer-reviewed academic sources and high-quality gray literature.}
    \label{fig:prisma}
\end{figure}

\subsection{Data Extraction}

We extracted structured data including:
\begin{itemize}
    \item Source metadata (type, date, authors)
    \item Protocol analyzed (EigenLayer, Symbiotic, Karak, Babylon, other)
    \item Threat categories identified
    \item Vulnerability types
    \item Mitigation strategies proposed
    \item Risk allocation mechanisms
    \item Research gaps identified
\end{itemize}

\subsection{Quality Assessment}

We assessed source quality using criteria adapted for mixed academic/gray literature:
\begin{itemize}
    \item \textbf{Academic Papers:} Peer review status, venue quality, methodology rigor
    \item \textbf{Security Audits:} Auditor reputation, scope, depth of analysis
    \item \textbf{Technical Analyses:} Author expertise, technical depth, evidence quality
    \item \textbf{Protocol Documentation:} Official status, technical completeness, version currency
\end{itemize}

\vspace{-0.3cm}
\section{Restaking Protocol Architectures}
\label{sec:architecture}

This section analyzes architectural designs of major restaking protocols, addressing RQ1.

\subsection{Architectural Components}

Restaking protocols share common components with distinct implementations:

\subsubsection{Staking Mechanisms}

\textbf{EigenLayer:} Supports native ETH restaking (via withdrawal credentials) and LST restaking (stETH, rETH, cbETH)~\cite{eigenlayer2024,eigenlayerwhitepaper2023}. Restakers can delegate to operators or run operators themselves. The protocol has undergone multiple security audits~\cite{certora2025,code4rena2023} and formal verification.

\textbf{Symbiotic:} Universal staking supporting diverse ERC-20 tokens~\cite{symbiotic2024,symbioticwhitepaper2024}, enabling broader participation beyond ETH ecosystem. Received \$29M Series A funding in 2025.

\textbf{Karak:} Multi-asset support including ETH, LSTs, LRTs, WBTC, stablecoins~\cite{karak2024}, expanding capital sources. Operates across multiple blockchains including Ethereum, Arbitrum, BSC, Blast, Mantle, and K2.

\textbf{Babylon:} Bitcoin-native staking through self-custodial vaults~\cite{babylon2024}, no wrapping required. Extends Bitcoin utility to secure other blockchains.

\subsubsection{Operator Architecture}

Operators run AVS software and receive delegations. Key design choices:
\begin{itemize}
    \item \textbf{Operator Requirements:} Technical capabilities, infrastructure, reputation
    \item \textbf{Delegation Model:} Direct delegation vs. liquid restaking tokens (LRTs)
    \item \textbf{Operator Economics:} Fee structures, reward distribution
    \item \textbf{Operator Slashing:} Conditions and penalties for misbehavior
\end{itemize}

\subsubsection{AVS Integration}

AVSs integrate with restaking protocols to leverage shared security:
\begin{itemize}
    \item \textbf{AVS Registration:} Onboarding process and requirements
    \item \textbf{Slashing Conditions:} Defined misbehaviors triggering penalties~\cite{certora2025,durvasula2024robust}
    \item \textbf{Operator Selection:} How AVSs choose operators
    \item \textbf{Upgrade Mechanisms:} Governance for AVS changes
    \item \textbf{AVS Examples:} Data availability layers, oracle networks, bridges~\cite{zhang2024proof,das2025reliablocks}
\end{itemize}

\subsubsection{Slashing Mechanisms}

Slashing penalizes operators/validators for misbehavior:
\begin{itemize}
    \item \textbf{Dual-Layer Slashing:} Penalties from both Ethereum and AVSs
    \item \textbf{Slashing Conditions:} Defined by each AVS
    \item \textbf{Slashing Caps:} Maximum penalties per validator
    \item \textbf{Appeal Mechanisms:} Dispute resolution processes
\end{itemize}

\subsection{Architectural Comparison}

Table~\ref{tab:architecture_comparison} compares architectural features across protocols. Figure~\ref{fig:protocol_timeline} illustrates the timeline of protocol launches and TVL growth.

\begin{table}[!ht]
    \centering
    \scriptsize
    \caption{Architectural comparison of major restaking protocols}
    \label{tab:architecture_comparison}
    \begin{tabular}{lcccc}
        \toprule
        \textbf{Feature} & \textbf{EigenLayer} & \textbf{Symbiotic} & \textbf{Karak} & \textbf{Babylon} \\
        \midrule
        Assets & ETH, LSTs & ERC-20 & Multi-asset & BTC \\
        Launch & Mid-2023 & 2024 & 2024 & 2024 \\
        TVL (2024) & \$20B+ & \$500M+ & \$300M+ & \$100M+ \\
        Cross-Chain & No & No & Yes & No \\
        Operator & Yes & Yes & Yes & Yes \\
        LRT & Yes & Yes & Yes & No \\
        Slashing Caps & Yes & TBD & TBD & TBD \\
        \bottomrule
    \end{tabular}
\end{table}

\begin{figure}[!ht]
    \centering
    \includegraphics[width=0.9\textwidth]{figures/protocol_timeline.pdf}
    \caption{Timeline of restaking protocol launches and TVL growth (2023-2024). EigenLayer launched in mid-2023 and reached \$20B+ TVL by June 2024. Symbiotic, Karak, and Babylon launched in 2024 with emerging adoption.}
    \label{fig:protocol_timeline}
\end{figure}

\vspace{-0.3cm}
\section{Threat Taxonomy}
\label{sec:threats}

This section presents our threat taxonomy addressing RQ2, classifying security threats into five primary categories.

\subsection{Category 1: Dual-Layer Slashing Risks}

Restaking introduces dual-layer slashing where validators face penalties from both Ethereum consensus and AVSs.

\textbf{Threat Description:}
\begin{itemize}
    \item Validator misbehavior can trigger slashing on Ethereum (e.g., double-signing, downtime)
    \item Simultaneously, AVS misbehavior triggers additional slashing
    \item Compounded losses exceed traditional staking penalties
    \item Cascading failures if slashing affects multiple AVSs
\end{itemize}

\textbf{Attack Scenarios:}
\begin{itemize}
    \item \textbf{Correlated Slashing:} Single validator error causes slashing across Ethereum and multiple AVSs
    \item \textbf{AVS Collusion:} Malicious AVSs coordinate slashing conditions to maximize validator penalties
    \item \textbf{Infrastructure Failures:} Network issues cause simultaneous downtime penalties
\end{itemize}

\textbf{Impact:} High - Can result in total loss of staked capital if slashing caps are insufficient. Durvasula and Roughgarden~\cite{durvasula2024robust} provide formal analysis showing that correlated slashing can significantly reduce the effective security buffer between attack costs and profits.

\textbf{Documented Incidents:}
While no major slashing events have occurred in restaking protocols as of 2024, the theoretical risk is substantial given the \$20B+ in restaked capital. Historical slashing events in Ethereum PoS provide context: validators have been slashed for double-signing and downtime, with penalties ranging from 0.5\% to 100\% of staked ETH depending on severity~\cite{ethereum2024staking}.

\subsection{Category 2: Smart Contract Vulnerabilities}

Restaking protocols involve complex smart contracts with multiple integration points. Smart contract security has been extensively studied~\cite{atzei2017survey,huang2021survey,chen2020defining,luu2016making}, but restaking introduces novel attack surfaces through protocol interactions.

\textbf{Threat Description:}
\begin{itemize}
    \item Bugs in restaking contracts (delegation, slashing, withdrawals)~\cite{certora2025,code4rena2023}
    \item Vulnerabilities in AVS contracts~\cite{zhang2024proof}
    \item Integration bugs between LST protocols and restaking~\cite{li2024mitigating}
    \item Reentrancy, access control, and logic errors~\cite{rodler2018sereum,zhang2020smart}
    \item Cross-protocol interaction vulnerabilities~\cite{chen2021sok}
\end{itemize}

\textbf{Attack Scenarios:}
\begin{itemize}
    \item \textbf{Withdrawal Exploits:} Bypassing withdrawal restrictions to drain funds
    \item \textbf{Slashing Bypass:} Avoiding legitimate slashing penalties
    \item \textbf{Delegation Manipulation:} Unauthorized delegation changes
    \item \textbf{LST Integration Bugs:} Exploiting interactions with LST protocols
\end{itemize}

\textbf{Documented Vulnerabilities:}
Security audits have identified vulnerabilities including:
\begin{itemize}
    \item \textbf{Code4rena (2023):} Identified two high-risk vulnerabilities: (1) H-01: Slot and block number proofs not required for withdrawal verification, potentially allowing multiple withdrawals, (2) H-02: Inability to slash queued withdrawals containing malicious strategies due to misplaced increment operation. Additionally identified multiple medium and low-risk findings.
    \item \textbf{Certora (2025):} Comprehensive formal verification of EigenLayer's slashing mechanism identified vulnerabilities in contract logic, state management, and access control. The audit covered core managers, storage, interfaces, and libraries.
    \item \textbf{Multiple Audits:} Additional security assessments have identified medium-risk findings in access control, state management, and integration points with LST protocols.
\end{itemize}

\textbf{Attack Surface Analysis:}
The smart contract attack surface includes:
\begin{itemize}
    \item \textbf{Restaking Contracts:} Delegation, withdrawal, slashing, and reward distribution logic
    \item \textbf{AVS Contracts:} Integration points, slashing conditions, operator management
    \item \textbf{LST Integration:} Interactions with stETH, rETH, cbETH protocols
    \item \textbf{LRT Protocols:} Liquid restaking token contracts and their interactions
\end{itemize}

\subsection{Category 3: Centralization Risks}

Economic incentives may concentrate validator power, creating systemic risks. Restaking protocols face inherent centralization pressures due to economies of scale in operator infrastructure and delegation patterns. Validator centralization in PoS systems has been studied extensively~\cite{chen2022rollup}, but restaking may exacerbate these trends.

\textbf{Threat Description:}
\begin{itemize}
    \item \textbf{Validator Concentration:} Large validators with superior infrastructure and reputation attract disproportionate delegations, dominating AVS security provision
    \item \textbf{Reduced Censorship Resistance:} Centralization reduces the number of independent validators, making censorship easier
    \item \textbf{Single Points of Failure:} If large validators are compromised (technically or through regulatory action), significant portions of restaked capital are at risk
    \item \textbf{Governance Capture:} Dominant operators can influence AVS governance decisions, including slashing conditions and protocol upgrades
    \item \textbf{Barrier to Entry:} High infrastructure and reputation requirements create barriers for new operators
\end{itemize}

\textbf{Attack Scenarios:}
\begin{itemize}
    \item \textbf{Validator Collusion:} Large operators coordinate to control AVS security, potentially extracting rents or censoring transactions
    \item \textbf{Governance Attacks:} Dominant validators influence AVS decisions to favor their interests, including modifying slashing conditions
    \item \textbf{Infrastructure Attacks:} Targeting centralized validator infrastructure (cloud providers, data centers) could compromise multiple large operators simultaneously
    \item \textbf{Regulatory Capture:} Regulatory action against large operators could disrupt significant portions of restaking security
\end{itemize}

\textbf{Empirical Evidence:}
Early data from EigenLayer shows concentration trends:
\begin{itemize}
    \item Top operators control significant portions of delegated capital
    \item Delegation patterns favor established validators with track records
    \item New operators face challenges attracting delegations
\end{itemize}

\textbf{Impact:} High - Centralization undermines core blockchain security properties and creates systemic risks.

\subsection{Category 4: Liquidity and Unstaking Challenges}

LST/LRT tokens introduce liquidity risks during market stress. Unlike traditional staking where assets are locked, restaking creates additional layers of liquidity dependencies.

\textbf{Threat Description:}
\begin{itemize}
    \item \textbf{LST Depegging:} Liquid staking tokens (stETH, rETH, cbETH) can depeg from ETH during market volatility, affecting restaking positions
    \item \textbf{Unstaking Delays:} Ethereum's unstaking queue (currently 4-5 days) prevents immediate exits, creating liquidity risks during stress events
    \item \textbf{LRT Liquidity Mismatches:} Liquid Restaking Tokens (LRTs) may face liquidity challenges, especially during market downturns or protocol issues
    \item \textbf{Cascading Liquidations:} If LST/LRT values drop significantly, leveraged positions may face liquidations, creating downward price pressure
    \item \textbf{Redemption Delays:} Converting LRTs back to underlying assets may face delays, especially during high unstaking demand
    \item \textbf{Protocol Interdependencies:} Restaking protocols depend on LST protocol stability; LST failures can cascade to restaking
\end{itemize}

\textbf{Attack Scenarios:}
\begin{itemize}
    \item \textbf{Depeg Attacks:} Attackers manipulate LST/LRT prices through market manipulation or protocol exploits, triggering liquidations and panic selling
    \item \textbf{Liquidity Crunch:} Mass unstaking requests during market stress overwhelm protocols, causing delays and potential losses
    \item \textbf{Arbitrage Exploits:} Exploiting price differences between LST/LRT and underlying assets, or between different LRT protocols
    \item \textbf{Flash Loan Attacks:} Using flash loans to manipulate LST/LRT prices or exploit redemption mechanisms
    \item \textbf{Protocol Run:} Coordinated withdrawals causing protocol insolvency or significant delays
\end{itemize}

\textbf{Historical Precedents:}
LST depegging events have occurred:
\begin{itemize}
    \item stETH depegging during Terra/Luna collapse (May 2022)
    \item Various LST depegging events during market stress
    \item These events demonstrate vulnerability of LST-based restaking
\end{itemize}

\textbf{Impact:} High - Liquidity risks can cause significant losses and protocol failures during market stress. Historical LST depegging events demonstrate the vulnerability: stETH depegged to 0.93 ETH during the Terra/Luna collapse in May 2022, causing significant losses for leveraged positions~\cite{trailofbits2023staking}.

Table~\ref{tab:lrt_comparison} compares major LRT protocols and their risk profiles.

\subsection{Category 5: Cryptoeconomic Attacks}

Game-theoretic vulnerabilities in restaking economics create opportunities for strategic manipulation and value extraction. Cryptoeconomic security in blockchain systems has been analyzed~\cite{zhou2020high,chen2021sok}, but restaking introduces multi-party games with complex incentive structures~\cite{durvasula2024robust}.

\textbf{Threat Description:}
\begin{itemize}
    \item \textbf{AVS Reward Manipulation:} AVSs may structure rewards to extract maximum value from operators while minimizing their own costs
    \item \textbf{Operator Fee Extraction:} Operators may charge excessive fees or structure fees to maximize their returns at delegator expense
    \item \textbf{Delegator Incentive Misalignment:} Delegators may not have sufficient information or incentives to make optimal operator selection decisions
    \item \textbf{Race Conditions:} Reward distribution mechanisms may have race conditions allowing front-running or manipulation
    \item \textbf{Free-Riding:} Operators may free-ride on others' security provision
    \item \textbf{Adverse Selection:} Low-quality operators may enter the market, diluting security
\end{itemize}

\textbf{Attack Scenarios:}
\begin{itemize}
    \item \textbf{AVS Collusion:} AVSs coordinate to structure slashing conditions and rewards to extract maximum value from operators while minimizing their own security costs
    \item \textbf{Operator Cartels:} Large operators collude to control AVS security, setting fees and conditions to maximize profits
    \item \textbf{MEV Extraction:} Validators extract maximal extractable value (MEV) at delegator expense, including transaction ordering and front-running
    \item \textbf{Reward Manipulation:} Operators manipulate reward distribution through strategic timing or protocol interactions
    \item \textbf{Information Asymmetry:} Operators with superior information exploit delegators who lack visibility into AVS risks and operator performance
    \item \textbf{Governance Manipulation:} Stakeholders manipulate governance to change slashing conditions or reward structures in their favor
\end{itemize}

\textbf{Game-Theoretic Analysis:}
Restaking creates complex multi-party games:
\begin{itemize}
    \item \textbf{Operator-AVS Game:} Operators choose AVSs to secure; AVSs set slashing conditions and rewards
    \item \textbf{Delegator-Operator Game:} Delegators choose operators; operators set fees and manage risk
    \item \textbf{Competition Among Operators:} Operators compete for delegations through fees, reputation, and performance
    \item \textbf{Competition Among AVSs:} AVSs compete for operator security through rewards and slashing conditions
\end{itemize}

\textbf{Impact:} Medium-High - Cryptoeconomic attacks can extract value and reduce protocol efficiency, though may not directly compromise security.

\begin{table}[!ht]
    \centering
    \scriptsize
    \caption{Comparison of major Liquid Restaking Token (LRT) protocols}
    \label{tab:lrt_comparison}
    \begin{tabular}{lcccp{2.5cm}}
        \toprule
        \textbf{Protocol} & \textbf{TVL} & \textbf{LSTs} & \textbf{Risk Model} & \textbf{Audit} \\
        \midrule
        Renzo & \$3B+ & stETH, rETH, cbETH & Diversified AVS & Yes \\
        EtherFi & \$4B+ & Native + LSTs & Native restaking & Yes \\
        Puffer & \$1.5B+ & Native ETH & Native restaking & Yes \\
        Kelp DAO & \$800M+ & stETH, rETH & Diversified & Yes \\
        \bottomrule
    \end{tabular}
\end{table}

\begin{table}[!ht]
    \centering
    \scriptsize
    \caption{Threat severity and frequency analysis}
    \label{tab:threat_severity}
    \begin{tabular}{lcccp{3cm}}
        \toprule
        \textbf{Threat} & \textbf{Severity} & \textbf{Likelihood} & \textbf{Impact} & \textbf{Documented} \\
        \midrule
        Dual-Layer Slashing & High & Medium & Total loss & Theoretical \\
        Smart Contract Bugs & High & Medium & \$100M+ & Audits found issues \\
        Centralization & Med-High & High & Governance & Early trends \\
        Liquidity Risks & High & Medium & Depegging & LST history \\
        Crypto Attacks & Medium & Low-Med & Value extract & Theoretical \\
        \bottomrule
    \end{tabular}
\end{table}

\vspace{-0.3cm}
\section{Mitigation Strategies}
\label{sec:mitigations}

This section synthesizes mitigation strategies addressing RQ3.

\subsection{Technical Mitigations}

Table~\ref{tab:mitigation_comparison} compares mitigation strategies across protocols.

\begin{table}[!ht]
    \centering
    \scriptsize
    \caption{Mitigation strategies comparison}
    \label{tab:mitigation_comparison}
    \begin{tabular}{lcccp{2.5cm}}
        \toprule
        \textbf{Mitigation} & \textbf{Eigen} & \textbf{Symb} & \textbf{Karak} & \textbf{Status} \\
        \midrule
        Slashing Caps & Yes & Planned & Planned & Active \\
        Formal Verif. & Yes & TBD & TBD & Certora 2025 \\
        Audits & Multiple & Limited & Limited & Ongoing \\
        Operator Req. & Yes & Yes & Yes & All \\
        Governance & On-chain & TBD & TBD & Active \\
        Diversification & Encouraged & Encouraged & Encouraged & User \\
        \bottomrule
    \end{tabular}
\end{table}

\subsubsection{Slashing Caps}

Limiting maximum slashing penalties per validator~\cite{durvasula2024robust}:
\begin{itemize}
    \item Prevents total capital loss from correlated failures
    \item Requires careful calibration to balance security and risk
    \item Implementation varies across protocols
    \item EigenLayer implements caps to limit maximum slashing per validator
    \item Theoretical analysis shows caps reduce effective security but protect against catastrophic losses
\end{itemize}

\subsubsection{Formal Verification}

Mathematical proofs of contract correctness~\cite{ahmed2021survey}:
\begin{itemize}
    \item Certora conducted comprehensive formal verification of EigenLayer slashing mechanism (February 2025)~\cite{certora2025}
    \item Verified slashing logic, state transitions, and access controls
    \item Reduces smart contract vulnerabilities significantly
    \item Limited adoption across restaking ecosystem (only EigenLayer has formal verification)
    \item Other protocols rely on traditional security audits
\end{itemize}

\subsubsection{Access Control}

Robust permission systems~\cite{rodler2018sereum,zhang2020smart}:
\begin{itemize}
    \item Multi-signature requirements for critical operations
    \item Time-locked upgrades
    \item Role-based access control
    \item Smart contract security patterns~\cite{luu2016making,atzei2017survey}
\end{itemize}

\subsection{Cryptoeconomic Mitigations}

\subsubsection{Risk Diversification}

Spreading restaking across multiple AVSs~\cite{durvasula2024robust}:
\begin{itemize}
    \item Reduces correlated failure risk
    \item Requires operator/delegator discipline
    \item May reduce rewards through diversification
    \item Theoretical analysis shows diversification improves robustness~\cite{durvasula2024robust}
    \item Practical challenges in implementation and monitoring
\end{itemize}

\subsubsection{Insurance Mechanisms}

Third-party insurance for slashing events:
\begin{itemize}
    \item Protects delegators from operator failures
    \item Market-based risk pricing
    \item Limited availability and coverage
    \item Similar to DeFi insurance mechanisms~\cite{chen2021sok}
\end{itemize}

\subsubsection{Governance Frameworks}

Transparent decision-making for AVS parameters:
\begin{itemize}
    \item Prevents governance capture~\cite{chen2022rollup}
    \item Enables community oversight
    \item Challenges in balancing efficiency and decentralization
    \item On-chain governance models~\cite{eigenlayer2024}
\end{itemize}

\vspace{-0.3cm}
\section{Risk Distribution Analysis}
\label{sec:risk_distribution}

This section analyzes risk allocation across stakeholders, addressing RQ4. Table~\ref{tab:risk_allocation} summarizes risk distribution across different stakeholder groups.

\begin{table}[!ht]
    \centering
    \scriptsize
    \caption{Risk allocation across stakeholders}
    \label{tab:risk_allocation}
    \begin{tabular}{lcccp{2.5cm}}
        \toprule
        \textbf{Stakeholder} & \textbf{Risks} & \textbf{Control} & \textbf{Impact} & \textbf{Mitigation} \\
        \midrule
        Operators & Dual slash, infra & High & High & Diversify, insure \\
        Delegators & Operator fail, depeg & Low-Med & High & Select ops, LRT \\
        AVSs & Low security, misbehave & Medium & Medium & Select ops, gov \\
        LST Issuers & Restaking risk & Medium & Medium & Limits, monitor \\
        End Users & AVS failures & Low & Low-Med & Select AVS \\
        \bottomrule
    \end{tabular}
\end{table}

\subsection{Stakeholder Risk Profiles}

\subsubsection{Validators/Operators}

\textbf{Risks:}
\begin{itemize}
    \item Dual-layer slashing penalties
    \item Infrastructure failures
    \item AVS misbehavior penalties
    \item Reputation damage
\end{itemize}

\textbf{Controls:}
\begin{itemize}
    \item Operator selection and monitoring
    \item Infrastructure redundancy
    \item Risk diversification across AVSs
\end{itemize}

\subsubsection{Delegators}

\textbf{Risks:}
\begin{itemize}
    \item Operator slashing affecting delegations
    \item LST/LRT depegging
    \item Unstaking delays
    \item Reward volatility
\end{itemize}

\textbf{Controls:}
\begin{itemize}
    \item Operator due diligence
    \item Diversification across operators
    \item Insurance products
    \item LRT liquidity management
\end{itemize}

\subsubsection{AVSs}

\textbf{Risks:}
\begin{itemize}
    \item Operator misbehavior
    \item Insufficient security (low operator participation)
    \item Slashing condition disputes
    \item Protocol upgrades
\end{itemize}

\textbf{Controls:}
\begin{itemize}
    \item Operator reputation systems
    \item Slashing condition clarity
    \item Governance mechanisms
\end{itemize}

\subsubsection{LST Issuers}

\textbf{Risks:}
\begin{itemize}
    \item Restaking affecting LST security
    \item Depegging from restaking events
    \item Regulatory concerns
\end{itemize}

\textbf{Controls:}
\begin{itemize}
    \item Restaking limits
    \item Risk monitoring
    \item Protocol integration controls
\end{itemize}

\vspace{-0.3cm}
\section{Comparative Protocol Analysis}
\label{sec:comparison}

This section provides detailed comparison of major restaking protocols.

\subsection{Security Model Comparison}

Table~\ref{tab:security_comparison} compares security features. Figure~\ref{fig:security_landscape} visualizes the security landscape across different threat categories.

\begin{table}[!ht]
    \centering
    \scriptsize
    \caption{Security feature comparison}
    \label{tab:security_comparison}
    \begin{tabular}{lcccc}
        \toprule
        \textbf{Feature} & \textbf{Eigen} & \textbf{Symb} & \textbf{Karak} & \textbf{Babylon} \\
        \midrule
        Slashing Caps & Yes & TBD & TBD & TBD \\
        Formal Verif. & Yes & TBD & TBD & TBD \\
        Audits & Multiple & Limited & Limited & Limited \\
        Operator Req. & Yes & Yes & Yes & Yes \\
        Governance & On-chain & TBD & TBD & TBD \\
        \bottomrule
    \end{tabular}
\end{table}

\begin{figure}[!ht]
    \centering
    \includegraphics[width=0.9\textwidth]{figures/security_landscape.pdf}
    \caption{Security landscape visualization showing threat categories, severity levels, and mitigation coverage across restaking protocols. The figure illustrates the relative security posture of different protocols and identifies areas requiring additional security measures.}
    \label{fig:security_landscape}
\end{figure}

\subsection{Risk Profile Comparison}

\textbf{EigenLayer:} Mature protocol with extensive audits, but faces centralization risks from dominance.

\textbf{Symbiotic:} Broader asset support increases attack surface but reduces ETH concentration risk.

\textbf{Karak:} Cross-chain support adds complexity and potential failure points.

\textbf{Babylon:} Bitcoin-native approach avoids wrapping risks but faces integration challenges.

\vspace{-0.3cm}
\section{Research Gaps and Future Directions}
\label{sec:gaps}

This section identifies critical research gaps addressing RQ5.

\subsection{Formal Security Models}

\textbf{Gap:} Lack of formal game-theoretic models analyzing restaking security. While Durvasula and Roughgarden~\cite{durvasula2024robust} provide initial analysis, comprehensive models are needed.

\textbf{Research Needs:}
\begin{itemize}
    \item Nash equilibrium analysis of operator/AVS interactions~\cite{durvasula2024robust}
    \item Formal proofs of security properties under restaking~\cite{ahmed2021survey,perez2021security}
    \item Slashing mechanism game theory and optimal design
    \item Risk model validation with empirical data
    \item Extension of PoS security models~\cite{kiayias2017ouroboros,buterin2017casper,buchman2016tendermint,kwon2014tendermint} to restaking contexts
    \item Formal verification of AVS designs~\cite{zhang2024proof}
\end{itemize}

\subsection{Empirical Incident Analysis}

\textbf{Gap:} Limited empirical analysis of restaking security incidents. While Ethereum slashing events have been documented~\cite{ethereum2024staking}, restaking-specific incidents are rare due to the field's novelty.

\textbf{Research Needs:}
\begin{itemize}
    \item Historical slashing event analysis~\cite{ethereum2024staking}
    \item LST/LRT depegging case studies~\cite{trailofbits2023staking}
    \item Operator failure post-mortems
    \item Correlation analysis of failures across AVSs
    \item Comparative analysis with traditional staking incidents
\end{itemize}

\subsection{Composability Risks}

\textbf{Gap:} Underexplored risks from restaking integration with DeFi. DeFi composability risks have been studied~\cite{chen2021sok,gao2021flashloan,zhou2020high}, but restaking adds new dimensions through LRT protocols~\cite{renzo2024,etherfi2024,puffer2024,kelp2024}.

\textbf{Research Needs:}
\begin{itemize}
    \item Restaking-DeFi protocol interactions and attack vectors~\cite{chen2021sok}
    \item LRT integration risks in lending protocols~\cite{renzo2024,etherfi2024}
    \item Cascading failure analysis across restaking, LST, and DeFi protocols
    \item Systemic risk assessment of interconnected protocols
    \item Flash loan attacks on restaking positions~\cite{gao2021flashloan,zhou2020high}
    \item Cross-chain restaking risks~\cite{karak2024,li2023survey,belchior2021survey}
\end{itemize}

\subsection{Long-Term Sustainability}

\textbf{Gap:} Limited analysis of long-term restaking security sustainability.

\textbf{Research Needs:}
\begin{itemize}
    \item Economic security evolution and long-term sustainability
    \item Centralization trend analysis~\cite{chen2022rollup}
    \item Regulatory impact assessment
    \item Protocol upgrade security~\cite{perez2021security}
    \item Interoperability with other blockchain systems~\cite{herlihy2018atomic,zhang2021cross,li2023survey}
\end{itemize}

\vspace{-0.3cm}
\section{Discussion}
\label{sec:discussion}

\subsection{Key Findings}

Our review synthesizes available literature on restaking protocol security and reveals several critical findings:

\begin{itemize}
    \item \textbf{Dual-Layer Slashing:} Represents the most significant novel risk~\cite{durvasula2024robust}, requiring careful slashing cap design and risk management. Theoretical analysis shows correlated slashing can reduce security buffers significantly.
    
    \item \textbf{Smart Contract Maturity:} Early protocols show vulnerabilities~\cite{certora2025,code4rena2023}, but security audits and formal verification are improving. EigenLayer's formal verification by Certora~\cite{certora2025} represents progress, but broader adoption is needed.
    
    \item \textbf{Centralization Pressures:} Economic incentives favor large operators, requiring governance mechanisms to maintain decentralization~\cite{chen2022rollup}. Early data shows concentration trends that need monitoring.
    
    \item \textbf{Liquidity Risks:} LST/LRT tokens introduce novel liquidity challenges~\cite{trailofbits2023staking} requiring careful protocol design. Historical depegging events demonstrate vulnerability.
    
    \item \textbf{Research Maturity:} Field is very new with limited academic treatment~\cite{durvasula2024robust,zhang2024proof,li2024mitigating}; significant research gaps exist in formal models, empirical analysis, and composability risks.
    
    \item \textbf{Protocol Diversity:} Four major protocols (EigenLayer, Symbiotic, Karak, Babylon) show different architectural choices and risk profiles, enabling comparative analysis.
\end{itemize}

\subsection{Implications for Practitioners}

\textbf{Protocol Designers:}
\begin{itemize}
    \item Implement slashing caps with careful calibration~\cite{durvasula2024robust,certora2025}
    \item Conduct comprehensive security audits~\cite{certora2025,code4rena2023,trailofbits2024,openzeppelin2024} and formal verification~\cite{ahmed2021survey,certora2025}
    \item Design governance to prevent centralization~\cite{chen2022rollup}
    \item Plan for liquidity stress scenarios~\cite{trailofbits2023staking}
    \item Follow smart contract security best practices~\cite{huang2021survey,chen2020defining,luu2016making,rodler2018sereum}
\end{itemize}

\textbf{Operators:}
\begin{itemize}
    \item Diversify across AVSs to reduce correlated risk~\cite{durvasula2024robust}
    \item Maintain robust infrastructure
    \item Monitor slashing conditions carefully~\cite{certora2025}
    \item Consider insurance for large operations
    \item Understand validator economics and incentive structures
\end{itemize}

\textbf{Delegators:}
\begin{itemize}
    \item Conduct operator due diligence
    \item Diversify across operators
    \item Understand LST/LRT risks~\cite{renzo2024,etherfi2024,puffer2024,kelp2024,trailofbits2023staking}
    \item Monitor protocol developments
    \item Be aware of liquidity risks~\cite{li2021sovereign,trailofbits2023staking}
\end{itemize}

\subsection{Limitations}

\begin{itemize}
    \item \textbf{Rapidly Evolving Field:} Restaking is very new~\cite{eigenlayer2024}; findings may become outdated quickly as protocols evolve.
    \item \textbf{Limited Academic Literature:} Heavy reliance on gray literature (audits~\cite{certora2025,code4rena2023,trailofbits2024}, blogs~\cite{blocksec2024,cobo2024,okx2024}) due to the field's novelty.
    \item \textbf{Emerging Protocols:} Some protocols (Symbiotic~\cite{symbiotic2024}, Karak~\cite{karak2024}, Babylon~\cite{babylon2024}) have limited operational history.
    \item \textbf{Selection Bias:} May miss non-English sources or unpublished analyses.
    \item \textbf{Data Availability:} Limited public data on slashing incidents and operator performance~\cite{ethereum2024staking}.
\end{itemize}

\vspace{-0.3cm}
\section{Conclusion}
\label{sec:conclusion}

Restaking protocols represent a transformative innovation in blockchain security~\cite{eigenlayer2024,symbiotic2024,karak2024,babylon2024}, enabling shared security models~\cite{wood2021polkadot} that enhance capital efficiency. However, our review reveals significant security challenges including dual-layer slashing risks~\cite{durvasula2024robust}, smart contract vulnerabilities~\cite{certora2025,code4rena2023}, centralization pressures, liquidity challenges~\cite{trailofbits2023staking}, and cryptoeconomic attacks~\cite{zhou2020high,chen2021sok}.

Key contributions of this review include: (1) comprehensive threat taxonomy across five categories, (2) comparative analysis of four major protocols, (3) synthesis of mitigation strategies, (4) risk distribution analysis, and (5) identification of critical research gaps.

Critical research gaps include formal game-theoretic models, empirical incident analysis, composability risk assessment, and long-term sustainability analysis. Addressing these gaps is essential for the secure evolution of the restaking ecosystem.

As restaking protocols continue to grow in adoption and complexity, comprehensive security analysis becomes increasingly critical. This review provides a foundation for future research and actionable guidelines for protocol designers, operators, and delegators.

\section*{Acknowledgments}

We thank the security research community, audit firms, and protocol teams for maintaining public documentation and analyses that enabled this review. We acknowledge the contributions of Certora, Code4rena, and other security auditors whose reports provided critical insights into restaking protocol vulnerabilities.

\bibliographystyle{splncs04}
\bibliography{restaking_references}

\end{document}
